\begin{abstract}
% No single large language model (LLM) is best for every query and every budget, making routing, i.e., deciding which model should answer each incoming query, a central mechanism for cost-effective LLM deployment. 
No single large language model (LLM) is optimal across all queries and budget constraints, making model routing essential for cost-effective LLM deployment.
% Existing routers span binary quality predictors, cost-aware cascades, graph-based routers, and agentic routers trained with reinforcement learning, yet their diverse formalisms and incompatible implementations make them difficult to compare, extend, and the evaluation of these routers lack a standard pipeline. 
Existing routers span binary quality predictors, cost-aware cascades, graph-based routers, and agentic routers, yet their diverse formalisms and incompatible implementations, coupled with the absence of a standardized evaluation pipeline, hinder fair comparison and further extension.
% In this paper, we first present a unified formulation of LLM routing as a sequential decision process, under which a router is fully specified by five components, including a context encoder, a model encoder, a scoring function, a decision rule, and a learning signal, and existing methods reduce to three families of single-turn, multi-turn, and personalized routing. 
In this paper, we present a unified formulation of LLM routing as a sequential decision process. Under this formulation, a router can be characterized in terms of five types of components: context encoders, model encoders, scoring functions, decision rules, and learning signals. Existing methods can then be organized into three families of single-turn, multi-turn, and personalized routing.
% Building on the formulation, we develop an automatic pipeline that constructs routing supervision by sweeping a candidate pool over benchmarks and evaluates any router on both response quality and inference cost under one protocol, yielding xRouteBench, a benchmark spanning general NLP, memory-augmented, multimodal, video, time-series, and personalized scenarios. 
Building on this formulation, we develop an automated pipeline that constructs routing supervision by systematically running a pool of candidate models across benchmarks and evaluates routers in terms of both response quality and inference cost under a unified protocol. The resulting benchmark, \textbf{xRouteBench}, spans generic LLM tasks, memory-augmented, vision (image and video), time-series, and personalized routing scenarios.
% Grounded in the formulation and pipeline, we present \textbf{\method}, an open-source library for standardized, modular implementation of LLM routers, where a user adds a router by implementing a single routing method and a loss, and more than 16 representative routers ship across the three families. 
Grounded in the formulation and pipeline, we present \textbf{\method}, an open-source infrastructure for standardized and modular implementation of LLM routers, where users can add a new router by implementing only a routing method and a loss function and access built-in implementations of more than 16 representative routers spanning all three families.
% With the library and benchmark, we conduct a systematic empirical study of LLM routing, showing that learned routing surpasses the strongest fixed model choice by 14.6\% relative, that cost pressure reverses the router ranking in favor of lightweight designs, and that user-conditioned routing yields consistent personalization gains.
Using the library and benchmark, we conduct a systematic empirical study of LLM routing and find that learned routers achieve a 14.6\% relative improvement over the strongest fixed-model baseline, router rankings reverse in favor of lightweight designs under tighter cost constraints, and user-conditioned routing delivers consistent personalization gains.
\end{abstract}
